% =========================================================
%  Capítulo 5 — Resultados
% =========================================================

\chapter{Resultados}
\label{cap:resultados}

Este capítulo apresenta os resultados obtidos com o desenvolvimento da plataforma,
abrangendo os testes realizados, as funcionalidades entregues e uma avaliação
qualitativa do sistema.

% -----------------------------------------------------------
\section{Funcionalidades Implementadas}
\label{sec:funcionalidades}

Ao final do desenvolvimento, as seguintes funcionalidades foram entregues:

\begin{itemize}
  \item \textbf{Autenticação e autorização:} cadastro de usuários, \textit{login}
        com JWT, redefinição de senha por e-mail e controle de acesso baseado em
        \textit{roles};

  \item \textbf{Organização acadêmica:} gerenciamento completo de períodos letivos,
        disciplinas, horários de aula, avaliações e registros de falta;

  \item \textbf{Atividades com \textit{checklist}:} criação e conclusão de
        atividades com lista de tarefas embutida, serializada como JSON;

  \item \textbf{Sessões de foco:} registro de sessões de estudo baseadas no método
        Pomodoro, com concessão de XP proporcional ao tempo dedicado;

  \item \textbf{Notas e materiais:} criação de notas de estudo e upload de
        materiais (arquivos) com armazenamento no AWS S3;

  \item \textbf{Planejador e lembretes:} criação de eventos e lembretes pessoais;

  \item \textbf{Gamificação:} sistema de XP, níveis, moedas e \textit{streak} com
        progressão baseada em ações do usuário;

  \item \textbf{Notificações:} alertas automáticos para atividades próximas do
        prazo, limite de faltas e subida de nível, com deduplicação.
\end{itemize}

% -----------------------------------------------------------
\section{Testes Realizados}
\label{sec:testes}

\subsection{Cobertura dos Testes Unitários}
\label{subsec:cobertura-testes}

[Apresente aqui os resultados dos testes unitários. Inclua uma tabela com as
classes testadas, número de casos de teste e a cobertura obtida. Execute
\texttt{./mvnw test} e insira o relatório.]

\begin{table}[H]
  \centering
  \caption{Resumo dos testes unitários implementados}
  \label{qdr:testes}
  \begin{tabular}{llc}
    \toprule
    \textbf{Classe testada}          & \textbf{Tipo}      & \textbf{N.º de casos} \\
    \midrule
    \texttt{UserTest}                & Modelo de domínio  & [N]                   \\
    \texttt{UserServiceImplTest}     & Serviço            & [N]                   \\
    \texttt{NotificationServiceImplTest} & Serviço        & [N]                   \\
    \texttt{UserControllerTest}      & Controlador        & [N]                   \\
    [Demais classes\ldots]           & --                 & --                    \\
    \bottomrule
  \end{tabular}
  \fonte{Elaborado pelo autor (\imprimirdata).}
\end{table}

\subsection{Casos de Teste Relevantes}
\label{subsec:casos-teste}

[Descreva os casos de teste mais importantes, especialmente os que validam
regras de negócio críticas como o sistema de XP/nível e a deduplicação de
notificações.]

% -----------------------------------------------------------
\section{Avaliação da Arquitetura}
\label{sec:avaliacao-arquitetura}

A adoção da Arquitetura Hexagonal trouxe benefícios observáveis durante o
desenvolvimento:

\begin{itemize}
  \item \textbf{Testabilidade:} os serviços de domínio podem ser testados de
        forma completamente isolada, sem necessidade de instanciar \textit{beans}
        do Spring, banco de dados ou outros componentes de infraestrutura;

  \item \textbf{Substituição de adapters:} a implementação de armazenamento
        pode ser trocada de AWS S3 para outro provedor sem alteração da lógica
        de domínio, bastando implementar a interface \texttt{FileStorageServicePort};

  \item \textbf{Legibilidade:} a separação clara de responsabilidades facilita
        a localização e manutenção do código, especialmente para novos
        contribuidores.
\end{itemize}

Como contrapartida, a arquitetura exige a criação de um número maior de classes
(duas camadas de mapeamento, interfaces de porta, adaptadores), o que aumenta
o \textit{boilerplate} inicial.

% -----------------------------------------------------------
\section{Discussão}
\label{sec:discussao}

[Discuta os resultados obtidos em relação aos objetivos definidos no
Capítulo~\ref{cap:introducao}. Analise pontos fortes e limitações da
solução. Relacione com o referencial teórico apresentado no
Capítulo~\ref{cap:referencial}.]
